\begin{textsample}{POS Dim 4 – 2015\_c6 – Score 14.00 – 2015\_c6/enade\_2015\_tec\_gest\_publ\_0d4a20dfe1.txt}  \label{ex:f4_pos_022}
CAPÍTULO 3 ANÁLISE TÉCNICA DA PROVA Este capítulo tem por objetivo apresentar o desempenho dos estudantes concluintes de Tecnologia em Gestão Pública no Enade/2015.
Para isso, foram calculadas as estatísticas básicas da prova em seu todo ( seção 3.1.1 ), bem como as estatísticas relacionadas aos Componentes de Formação Geral ( seção 3.1.2 ) e de Conhecimento Específico da Área ( seção 3.1.3 ).
Dadas as suas características, foram analisadas, em separado, as questões objetivas ( seção 3.2 ) e as questões discursivas ( seção 3.3 ).
Para as questões objetivas, foram disponibilizados os índices de facilidade e de discriminação bisserial, também em separado, para os Componentes de Formação Geral ( seção 3.2.1 ) e de Conhecimento Específico ( seção 3.2.2 ).
De cada componente, uma das questões foi escolhida para exemplificar a análise gráfica, relacionando as alternativas escolhidas pelos estudantes ( inclusive o gabarito ) com o número de acertos no componente.
O Anexo I apresenta a íntegra da análise gráfica para todas as questões objetivas.
Para cada uma das questões discursivas, os conteúdos dos tipos mais comuns de respostas dos estudantes são apresentados e comparados com o padrão de respostas esperado ( ver Anexo VIII com o padrão de respostas ).
Tomando como base as duas questões discursivas do Componente de Formação Geral, a seção 3.3.1.6 apresenta comentários sobre a correção das respostas com respeito à Língua Portuguesa.
Nas Tabelas 3.1 e 3.3 são apresentados o tamanho da população inscrita, ausente, presente e percentual de faltosos, respectivamente, por Grande Região e por Categoria Administrativa e Organização Acadêmica da IES.
As demais tabelas apresentam as seguintes estatísticas das notas15 : média do desempenho na prova, erro padrão da média, desvio padrão, nota mínima, mediana e nota máxima para cada um de seus componentes.
Tais estatísticas contemplam o total de estudantes concluintes da Área de Tecnologia em Gestão Pública inscritos e presentes à prova do Enade/2015, tendo em vista agregações ou por Grandes Regiões e o país como um todo, ou por Categoria Administrativa e Organização Acadêmica da IES.
Em relação aos gráficos de distribuição de notas, o intervalo considerado foi de 10 unidades, aberto à esquerda e fechado à direita, com exceção do primeiro intervalo, [ 0 ; 10 ], fechado em ambos os extremos.
Para os gráficos de distribuição das notas das 15 Uma definição dessas estatísticas pode ser encontrada no Glossário. questão em branco16 e nota zero. 3.1 CONDIÇÃO DE PRESENÇA E ESTATÍSTICAS BÁSICAS DA PROVA Esta seção apresenta a condição de presença, além das estatísticas selecionadas e histogramas da nota geral ( 3.1.1 ) e de cada componente : Formação Geral ( 3.1.2 ) e Conhecimento Específico ( 3.1.3 ).
São também apresentadas estatísticas selecionadas de subpopulações caracterizadas por Grande Região, Categoria Administrativa e Organização Acadêmica. 3.1.1 Condição de Presença e Estatísticas Básicas Gerais A Tabela 3.1 apresenta a Condição de Presença na prova, por Grande Região, dos estudantes concluintes de Tecnologia em Gestão Pública.
A população total de Inscritos foi de 5.336.
Destes, 4.332 estiveram Presentes, sendo 18,8 \% o índice de não comparecimento.
A Região de maior abstenção foi a Centro-Oeste ( 26,5 \%, ) e a de menor abstenção foi a Norte ( 16,6 \% ).
Já a Tabela 3.2 apresenta as Estatísticas Básicas da prova, por Grande Região, dos estudantes concluintes de Tecnologia em Gestão Pública.
A Média das notas da prova como um todo ( nas seções seguintes serão analisados os Componentes de Formação Geral e de Conhecimento Específico ) foi 45,9, sendo que os alunos da região Sul obtiveram a Média mais baixa ( 44,5 ), e os da região Sudeste obtiveram a Média mais alta ( 51,1 ).
As demais médias foram : 50,2 na região Norte ; 48,9 na região Nordeste ; e 44,8 na região Centro-Oeste.
O Desvio padrão para o Brasil como um todo foi 15,0, sendo o maior Desvio padrão encontrado na região Centro-Oeste ( 15,7 ) e o menor, na região Nordeste ( 13,7 ), indicando uma dispersão um pouco menor das notas desta última região.
A região que obteve a maior nota Máxima foi a Sudeste ( 92,4 ), ao passo que a região que atingiu a menor nota Máxima foi a Nordeste ( 81,0 ).
A Mediana do Brasil como um todo foi 45,2, sendo a maior Mediana obtida na região Sudeste ( 50,6 ), e a menor obtida no Sul ( 43,8 ).
A nota Mínima foi zero apenas nas regiões Sudeste e Sul, enquanto nas demais regiões foi : Norte ( 8,9 ) ; Nordeste ( 13,8 ) ; e Centro-Oeste ( 5,1 ).
Considerando -se as notas segundo Grande Região, observa -se que existe diferença estatisticamente significativa ao nível de 95 \% entre as menores Média, obtidas nas regiões Sul ( 44,5 ) e Centro-Oeste ( 44,8 ), e as médias de todas as demais regiões.
O comportamento das notas dos estudantes de todo o Brasil pode ser observado no Gráfico 3.1 que apresenta um histograma com a distribuição das mesmas.
Essa é uma distribuição unimodal com a moda no intervalo ( 40 ; 50 ].
A Tabela 3.3 apresenta informações referentes à Condição de Presença desagregadas por Categoria Administrativa e por Organização Acadêmica.
Da população total de Inscritos, 4.365 são de IES Privada e 971 de IES Públicas.
Em relação à Organização Acadêmica, a maior participação foi obtida por estudantes de Centros Universitários ( 2.939 ), seguido de Universidades ( 2.023 ) e Faculdades ( 374 ).
A Categoria Administrativa de maior abstenção foi a Pública ( 19,6 \% ), e entre as Organizações Acadêmicas foi a das Faculdades ( 25,4 \% ), que ficou acima da abstenção nacional de 18,8 \%.
Já a Tabela 3.4 apresenta informações referentes à Média da nota final desagregadas por Categoria Administrativa e por Organização Acadêmica.
A Média das notas da prova como um todo foi 45,9.
Em relação à Categoria Administrativa, os estudantes das IES Públicas obtiveram Média mais alta ( 50,6 ), e os das IES Privadas obtiveram Média mais baixa ( 44,9 ), que a Média nacional.
Observa -se que existe diferença estatisticamente significativa entre as médias das notas das IES Públicas e Privadas.
A diferença entre as médias das regiões Sudeste e Sul ( 6,6 ), a maior e a menor Média, é superior à diferença entre IES Públicas e Privadas ( 5,7 ), caracterizando uma maior diversidade regional do que administrativa.
No tocante à Organização Acadêmica, apenas as Universidades obtiveram Média mais alta ( 49,8 ) que a nacional.
As demais médias foram : 44,1 nos Centros Universitários e 40,7 nas Faculdades.
O Desvio padrão para as IES Privadas ( 15,1 ) e para as Faculdades ( 15,8 ) foram superiores ao do Brasil como um todo ( 15,0 ), indicando uma dispersão um pouco maior das notas nesta Categoria Administrativa e nesta Organização Acadêmica.
Constata -se que existe diferença estatisticamente significativa ao nível de 95 \% nas médias das notas dos estudantes dos tipos de Organização Acadêmica. 3.1.2 Estatísticas Básicas no Componente de Formação Geral A Tabela 3.5 apresenta as Estatísticas Básicas em relação ao componente da prova que avalia a Formação Geral dos estudantes concluintes.
Os alunos de todo Brasil obtiveram desempenho médio de 53,6.
Quanto à variabilidade, o Desvio padrão das notas dos estudantes do Brasil, como um todo, foi 17,4.
A maior Média foi obtida na região Sudeste ( 58,4 ), e a menor, na região Centro-Oeste ( 52,1 ).
As demais médias foram : 58,2 na região Norte ; 55,5 na região Nordeste ; e 52,4 na região Sul.
Já o maior Desvio padrão foi obtido na região Centro-Oeste ( 17,9 ), e o menor, na região Norte ( 15,8 ).
Os demais desvios padrões foram : 16,3 na região Nordeste, 16,9 da região Sudeste, e 17,4 na região Sul.
A maior nota no Componente de Formação Geral da prova do Enade/2015 foi 97,6, obtida por, pelo menos, um aluno na região Nordeste.
A menor nota Máxima foi obtida na região Centro-Oeste ( 92,2 ).
A Mediana do Brasil, como um todo, foi 54,8, sendo a menor Mediana encontrada na região Sul ( 53,6 ) ; e a maior encontrada na região Sudeste ( 60,1 ).
A nota Mínima nesta parte foi zero em quase todas as regiões, exceto nas regiões Norte ( 15,0 ) e Nordeste ( 7,5 ).
Considerando -se as notas segundo Grande Região, observa -se que existe diferença estatisticamente significativa entre a maior Média das notas do Componente de Formação Geral, obtida na região Sudeste ( 58,4 ), e na região Centro-Oeste ( 52,1 ), a menor Média.
O Gráfico 3.2 propicia a avaliação do desempenho dos estudantes do Componente de Formação Geral, a partir do histograma da distribuição das notas correspondentes.
A distribuição é unimodal, com moda em ( 50 ; 60 ], enquanto na prova, como um todo ( Gráfico 3.1 ), a moda foi alcançada um intervalo abaixo ( 40 ; 50 ].
Nota -se, ainda que, no Gráfico 3.2, as notas apresentam maior dispersão do que no Gráfico 3.1 ( distribuição das notas da prova ), confirmado pela comparação dos desvios padrões : 15,0 para a nota da prova como um todo e 17,4 para o Componente de Formação Geral.
Na Tabela 3.6 são apresentadas as informações referentes ao desempenho dos concluintes do Componente de Formação Geral, em diferentes agregações : Categoria Administrativa e Organização Acadêmica.
Observa -se que existe diferença estatisticamente significativa entre as médias dos tipos de Categoria Administrativa.
A maior Média foi obtida por estudantes de IES Públicas ( 57,1 ), com uma diferença estatisticamente significativa da obtida por estudantes de IES Privadas ( 52,8 ).
Considerando -se o tipo de Organização Acadêmica, nota -se que há diferença estatisticamente significativa entre as médias dos tipos de Organização Acadêmica. 3.1.3 Estatísticas Básicas do Componente de Conhecimento Específico A Tabela 3.7 apresenta as Estatísticas Básicas referentes ao Componente de Conhecimento Específico da área de Tecnologia em Gestão Pública.
A Média do desempenho dos alunos do Brasil, como um todo, foi 43,3.
A maior Média foi obtida na região Sudeste ( 48,6 ), e a menor, na região Sul ( 41,9 ).
As demais médias foram : 47,5 na região Norte ; 46,6 na região Nordeste ; e 42,4 na região Centro-Oeste.
Quanto à variabilidade das notas, o Desvio padrão do Brasil, como um todo, foi 16,2, sendo o maior Desvio padrão observado na região Centro-Oeste ( 17,1 ), e o menor, na região Nordeste ( 15,2 ).
Os demais desvios foram : 15,3 na região Norte, 16,3 na região Sudeste e 15,9 na região Sul.
A Mediana das notas dos estudantes de todo o Brasil foi 42,3.
A maior Mediana ocorreu na região Sudeste ( 48,2 ), e a menor, na região Sul ( 40,5 ).
As demais medianas foram : 48,0 na região Norte ; 45,4 na região Nordeste ; e 40,7 na região Centro-Oeste.
A nota Máxima do Brasil, como um todo, foi 93,2, sendo obtida por, pelo menos, um aluno na região Sudeste.
As demais notas máximas foram : 80,7 na região Norte ; 80,5 na região Nordeste ; 91,2 na região Sul ; e 84,9 na região Centro-Oeste.
A nota Mínima foi zero apenas nas regiões Sudeste e Sul, enquanto nas demais regiões foi : 6,8 nas regiões Norte e Centro-Oeste ; e 8,6 na região Nordeste.
Observa -se que existe diferença estatisticamente significativa da menor Média das notas do Componente de Conhecimento Específico da região Sul ( 41,9 ) em relação à maior Média da região Sudeste ( 48,6 ).
O intervalo de confiança da região Norte ( 5,3 ) é significativamente maior do que de todas as outras.
Assim como os Gráficos 3.1 e 3.2, o Gráfico 3.3, apresentado a seguir, permite uma avaliação do desempenho de concluintes de Tecnologia em Gestão Pública em relação ao Componente de Conhecimento Específico com um histograma da distribuição das notas correspondentes.
Dentre as três distribuições apresentadas, esta é a mais concentrada nas notas baixas.
Esta também é uma distribuição unimodal, e o grupo modal é o ( 30 ; 40 ], abaixo do grupo modal tanto da prova como um todo quanto para a Formação Geral.
A Tabela 3.8 apresenta uma comparação dos resultados em relação à Categoria Administrativa e à Organização Acadêmica, agora levando em conta o desempenho de alunos do Componente de Conhecimento Específico da prova.
No que se refere à Organização Acadêmica, a maior Média foi das Universidades ( 47,5 ), vindo a seguir a dos Centros Universitários ( 41,3 ) e, depois, a das Faculdades ( 38,3 ).
O maior Desvio padrão, e acima do valor para o Brasil, como um todo, foi o das Faculdades ( 16,8 ).
As Faculdades também obtiveram a maior nota Máxima ( 93,2 ) e a menor Mediana ( 35,9 ).
Por sua vez, as Universidades obtiveram 86,3 e 47,6 como nota Máxima e Mediana, respectivamente.
Já os Centros Universitários, obtiveram 91,2 e 39,9 como nota Máxima e Mediana, respectivamente.
A nota Mínima foi zero para duas das Organizações Acadêmicas, exceto para as Faculdades ( 6,8 ).
Quanto à Categoria Administrativa, observa -se um comportamento semelhante àquele da parte de Formação Geral e à prova como um todo, ou seja, existe diferença estatisticamente significativa entre as médias das IES Públicas ( 48,4 ) e IES Privadas ( 42,2 ).
Considerando -se o tipo de Organização Acadêmica, nota -se que há diferença estatisticamente significativa entre as médias dos tipos de Organização Acadêmica 3.2 ANÁLISE DAS QUESTÕES OBJETIVAS Esta seção apresenta estatísticas selecionadas e histogramas das Questões Objetivas de Formação Geral ( 3.2.1 ).
São também apresentadas e comparadas as médias das subpopulações caracterizadas por Grande Região. 3.2.1 Componente de Formação Geral A Tabela 3.9 apresenta as Estatísticas Básicas relativas às oito questões objetivas do componente da prova que abrange a Formação Geral dos estudantes.
A Média do Brasil foi 55,0.
A menor Média foi encontrada na região Centro-Oeste ( 52,9 ), e a maior, na região Sudeste ( 60,3 ).
As demais médias foram : 58,2 na região Norte ; 54,3 na região Nordeste ; e 54,1 na região Sul.
O Desvio padrão do Brasil foi 21,0, sendo o maior Desvio padrão encontrado na região Sudeste ( 21,8 ), e o menor, na região Norte ( 20,1 ).
Os demais desvios foram : 20,5 na região Nordeste ; 20,7 na região Sul ; e 21,1 na região Centro-Oeste.
As medianas ( 50,0 ), as notas máximas ( 100,0 ) e as notas mínimas ( 0,0 ) foram iguais para quase todas as regiões, as exceções foram a Mediana nas regiões Norte e Sudeste ( 62,5 cada ) e a nota Mínima na região Norte ( 12,5 ).
A Tabela 3.10 apresenta o Índice de Facilidade e o Índice de Discriminação ( ponto bisserial ) para cada uma das questões objetivas do Componente de Formação Geral.
Quanto ao Índice de Facilidade, foram usadas as seguintes cores para diferenciar o nível de dificuldade da questão :  Azul para as questões classificadas com índice Muito fácil ( \&gt ; =0,86 ), verde para as questões classificadas com índice Fácil ( 0,61 a 0,85 ), amarelo para as questões classificadas com Médio ( 0,41 a 0,60 ), vermelho para as questões classificadas com Difícil ( 0,16 a 0,40 ) e roxo para as questões classificadas com Muito difícil ( \&lt ; =0,15 ).
Já quanto ao Índice de Discriminação, foram usadas as seguintes cores para qualificar a questão :  As questões classificadas com índice Fraco receberam a cor vermelho ( \&lt ; =0,19 ), as classificadas com Médio receberam a cor amarelo ( 0,20 a 0,29 ), as classificadas com Bom receberam a cor verde ( 0,30 a 0,39 ) e as classificadas com Muito bom ( \&gt ; =0,40 ) receberam a cor azul.
As questões objetivas do Componente de Formação Geral, segundo o Índice de Facilidade, foram assim avaliadas : das oito questões, nenhuma teve o Índice de Facilidade classificado como Muito fácil.
Quatro questões foram tidas como Fácil, por terem índice de acertos situado na faixa entre 0,61 e 0,85 ( de 61,0 \% a 85,0 \% de acertos ).
Duas questões foram consideradas de dificuldade Médio, situando -se no intervalo entre 0,41 e 0,60 do Índice de Facilidade, ou seja, houve entre 41,0 \% e 60,0 \% de acertos, enquanto outras duas questões foram classificadas na categoria Difícil, situando -se no intervalo entre 0,16 e 0,40.
Por fim, nenhuma das questões apresentou menos de 15 \% de acertos, razão pela qual seria classificada como Muito difícil.
Como já comentado, para análise das questões objetivas, relativas à Formação Geral, segundo o poder de discriminação, utilizou -se o Índice de Discriminação ( ponto A Tabela 3.10 apresenta o Índice de Facilidade e o Índice de Discriminação ( ponto bisserial ) para cada uma das questões objetivas do Componente de Formação Geral.
Quanto ao Índice de Facilidade, foram usadas as seguintes cores para diferenciar o nível de dificuldade da questão :  Azul para as questões classificadas com índice Muito fácil ( \&gt ; =0,86 ), verde para as questões classificadas com índice Fácil ( 0,61 a 0,85 ), amarelo para as questões classificadas com Médio ( 0,41 a 0,60 ), vermelho para as questões classificadas com Difícil ( 0,16 a 0,40 ) e roxo para as questões classificadas com Muito difícil ( \&lt ; =0,15 ).
Já quanto ao Índice de Discriminação, foram usadas as seguintes cores para qualificar a questão :  As questões classificadas com índice Fraco receberam a cor vermelho ( \&lt ; =0,19 ), as classificadas com Médio receberam a cor amarelo ( 0,20 a 0,29 ), as classificadas com Bom receberam a cor verde ( 0,30 a 0,39 ) e as classificadas com Muito bom ( \&gt ; =0,40 ) receberam a cor azul.
As questões objetivas do Componente de Formação Geral, segundo o Índice de Facilidade, foram assim avaliadas : das oito questões, nenhuma teve o Índice de Facilidade classificado como Muito fácil.
Quatro questões foram tidas como Fácil, por terem índice de acertos situado na faixa entre 0,61 e 0,85 ( de 61,0 \% a 85,0 \% de acertos ).
Duas questões foram consideradas de dificuldade Médio, situando -se no intervalo entre 0,41 e 0,60 do Índice de Facilidade, ou seja, houve entre 41,0 \% e 60,0 \% de acertos, enquanto outras duas questões foram classificadas na categoria Difícil, situando -se no intervalo entre 0,16 e 0,40.
Por fim, nenhuma das questões apresentou menos de 15 \% de acertos, razão pela qual seria classificada como Muito difícil.
Como já comentado, para análise das questões objetivas, relativas à Formação Geral, segundo o poder de discriminação, utilizou -se o Índice de Discriminação ( ponto bisserial ).
Nesta análise, as questões foram assim avaliadas : seis das oito questões apresentaram índice acima ou igual a 0,40 e, assim, foram classificadas com índice Muito bom para esse grupo de estudantes.
Duas questões tiveram Índice de Discriminação Bom, com valor entre 0,30 e 0,39, para esse grupo de estudantes.
Nenhuma questão teve nível Médio ou Fraco de discriminação para esse grupo de estudantes.
O Índice de Facilidade variou de 0,31 a 0,70, e o de Discriminação, de 0,32 a 0,53.
As seis questões com Índice de Discriminação Muito bom figuraram entre os diversos níveis de dificuldade desse conjunto : quatro classificadas na categoria Fácil ( questões 3, 5, 6 e 8 ) do Índice de Facilidade, uma, na categoria Médio ( questão 2 ) e uma, na categoria Difícil ( questão 7 ).
Em particular, a questão 2 foi a que apresentou o maior poder discriminatório, com índice 0,53, porém foi considerada média em termos de facilidade, com uma proporção de 0,59 acertos.
O máximo de acertos foi alcançado pela questão 6 com um Índice de Facilidade de 0,70.
A questão de número 1 apresentou um Índice de Facilidade de 0,49, ou seja, quase a metade dos estudantes conseguiu resolvê -la, dentro do universo de participantes.
Seu Índice de Discriminação foi Bom ( 0,32 ).
Já a questão 7 obteve Índice de Discriminação Muito bom, 0,51, porém seu Índice de Facilidade foi Difícil ( 0,40 ).
O Gráfico 3.4, para exemplificar, analisa o comportamento da questão de número 2 de Formação Geral.
Trata -se de uma questão considerada mediana, com relação à facilidade e a que obteve o maior Índice de Discriminação dessa parte da prova.
Neste gráfico, cada uma das cinco curvas representa o percentual de respostas em determinada alternativa da questão, em função do número de acertos dos estudantes nessa parte da prova ( Formação Geral/Múltipla Escolha ), antes de possíveis eliminações pelo critério do ponto bisserial.
A curva em preto corresponde à alternativa C, a correta para esta questão.
Assim, observa -se que entre os estudantes com menor número de acertos, nessa parte do exame, a situação mais frequente foi a escolha da alternativa D ( em roxo ), incorreta.
Por exemplo, entre os estudantes que acertaram três questões, 32,1 \% escolheram a alternativa D, 33,2 \% a alternativa C ( correta ), 14,8 \% a alternativa A, 13,3 \% a alternativa B e 6,2 \% a alternativa E.
Entre os que acertaram três respostas entre as questões de múltipla escolha de formação geral, 0,1 \% deixou a questão em branco e 0,3 \% marcou mais de uma alternativa, invalidando a questão.
À medida em que o número de acertos aumenta, indicando desempenho melhor nessa parte da prova, aumenta concomitantemente a proporção de estudantes que selecionaram a alternativa correta C, atingindo 100 \% para os estudantes com 8 acertos.
Essa análise permite verificar como a questão discriminou os grupos de desempenho, justificando o alto índice obtido na questão.
Cumpre notar que não é possível inferir deste gráfico nem o índice de facilidade ( que seria uma média da proporção ponderada pela quantidade de alunos com cada uma das notas ), nem o índice de discriminação ponto bisserial, por razão equivalente.
No caso extremo no qual a grande concentração dos acertos dos alunos fosse abaixo de 4, o índice de facilidade seria obrigatoriamente abaixo de 40 \% ( neste exemplo ).
Caso a concentração fosse em 6 acertos ou mais, o índice seria obrigatoriamente acima de 60 \%.
Os gráficos relativos às demais questões de Formação Geral constam do Anexo I. 3.2.2 Componente de Conhecimento Específico A Tabela 3.11 apresenta as Estatísticas Básicas em relação às questões objetivas do Componente de Conhecimento Específico da prova de Tecnologia em Gestão Pública por Grande Região.
A Média do Brasil deste componente foi de 43,9.
A menor Média foi observada na região Sul ( 42,5 ), e a maior, na região Sudeste ( 49,1 ).
O Desvio padrão de todo o Brasil foi 16,7, sendo o menor Desvio padrão encontrado na região Nordeste ( 15,3 ), e o maior, na região Centro-Oeste ( 17,4 ).
A Mediana de todo o Brasil foi 44,0, a mesma encontrada nas regiões Nordeste e Centro-Oeste, enquanto nas demais regiões a Mediana foi : 48,0 na região Norte e Sudeste ; e 40,0 na região Sul.
A nota Máxima da prova foi 92,0, obtida nas questões objetivas do Componente de Conhecimento Específico, por, pelo menos, um aluno nas regiões Sudeste e Sul, e nas demais regiões, a nota Máxima da prova foi : 88,0 na região Norte ; 80,0 na região Nordeste ; e 84,0 na região Centro-Oeste.
A nota Mínima foi zero apenas nas regiões Sudeste e Sul, enquanto nas demais regiões foi : 8,0 nas regiões Norte e Nordeste ; e 4,0 na região Centro-Oeste.
A Tabela 3.12 apresenta os Índices de Facilidade e Discriminação ( ponto bisserial ) das questões objetivas do Componente de Conhecimento Específico para os estudantes de Tecnologia em Gestão Pública.
Para facilitar a diferenciação das questões, usaremos as mesmas cores da Tabela 3.10 para as diferentes classificações dos Índices de Facilidade e de Discriminação.
Dentre as questões objetivas da parte da prova relativa ao Componente de Conhecimento Específico, uma questão foi anulada pela Comissão, a questão 22.
Desse modo, a classificação, quanto ao Índice de Facilidade, foi estabelecida com base em 26 questões.
A partir dos índices obtidos, pode -se concluir que metade das questões objetivas da prova foi considerado, pelo menos, Difícil : das 26 questões válidas, 13 foram classificadas como Difícil ( a classificação modal ) ou como Muito difícil.
Cinco questões foram classificadas como Fácil, e outras oito, consideradas como Médio.
Nenhuma questão foi considerada Muito fácil.
Já quanto aos índices de discriminação das questões objetivas do Componente de Conhecimento Específico da prova, tem -se como resultado a seguinte classificação : 11 das 26 questões foram consideradas como boas, enquanto nove delas tiveram Índice de Discriminação Muito bom.
Assim, para a maioria das questões – 20 em 26 – os Índices de Discriminação foram Bom ou Muito bom.
Dentre as demais, cinco delas foram classificadas como Médio e uma outra como Fraco, sendo seis, por conseguinte, a quantidade de questões nos dois patamares mais baixos de discriminação.
Constata -se, assim, que a prova – no que se refere ao Componente de Conhecimento Específico – possuía boa capacidade de discriminar entre aqueles que dominam ou não o conteúdo.
Dentre as questões que alcançaram os maiores índices de discriminação, nove delas, as de números 12, 14, 16, 17, 18, 19, 30, 31 e 32, foram classificadas com Índice de Discriminação Muito bom, situando -se no intervalo de 0,40 a 0,54 do índice, quatro delas ( questões 16, 18, 19 e 30 ) foram classificadas na categoria Fácil, quanto ao Índice de Facilidade, três delas ( questões 12, 31 e 32 ), na categoria Médio, e outras duas ( questões 14 e 17 ), na categoria Difícil.
A questão de número 25 foi a mais difícil dentre as 26 questões específicas válidas, com baixo Índice de Facilidade, apenas 15,0 \% de acertos.
Essa questão apresentou poder discriminatório Médio, 0,25, o que comprova ter sido esta questão uma das mais difíceis para os estudantes.
Destaca -se, também, a questão 23, com Índice de Facilidade 0,24, o que, em termos percentuais, corresponde a 24,0 \% de estudantes que responderam acertadamente.
Já 0,13 foi o seu Índice de Discriminação, e, portanto, tal questão foi considerada inadequada pelo critério ponto bisserial e, por isso, foi eliminada do cômputo da nota final.
Nenhuma outra questão foi considerada inadequada pelo critério ponto bisserial.
A título de exemplo das análises do comportamento das questões objetivas, o Gráfico 3.5 analisa a questão 32 do Componente de Conhecimento Específico.
Apesar de esta questão ter sido considerada pelos estudantes avaliados como uma questão de classificação Médio como nível de facilidade da prova, apresentou Índice de Facilidade 0,50, ou seja, 50,0 \% dos estudantes assinalaram acertadamente a opção D, correspondente ao gabarito.
Seu Índice de Discriminação foi igual a 0,54, classificado como Muito bom, também sendo esta questão a que apresentou o maior índice discriminatório.
Neste gráfico, cada uma das cinco curvas representa o percentual de respostas em determinada alternativa da questão 32, em função do número de acertos dos estudantes nessa parte da prova, antes de possíveis eliminações de questões pelo critério do ponto bisserial.
A alternativa correta D, representada no gráfico pela curva em roxo, foi escolhida em maiores proporções pelos alunos com desempenho melhor nessa parte da prova.
Já as alternativas incorretas, também denominadas distratores, foram selecionadas, principalmente, por aqueles com notas mais baixas.
Neste caso também a soma não é sempre 100 \% por causa das questões não respondidas ou com mais de uma opção marcada.
Aqueles com nota zero, na sua quase totalidade deixaram esta questão em branco ou marcaram mais de uma alternativa, comportamento considerado inválido.
A proporção de alunos que selecionou a resposta correta D aumenta gradativamente, chegando a atingir 100 \% para 22 acertos ou mais, enquanto a proporção dos que escolheram alternativas incorretas decai, a partir do primeiro acerto, como função do número de acertos nessa parte da prova.
Os gráficos relativos às demais questões do Conhecimento Específico constam do Anexo I. 3.3 ANÁLISE DAS QUESTÕES DISCURSIVAS Esta seção apresenta estatísticas selecionadas e histogramas das Questões Discursivas de Formação Geral ( 3.3.1 ) e Conhecimento Específico ( 3.3.2 ).
São também apresentadas e comparadas as médias de subpopulações caracterizadas por Grande Região. 3.3.1 Componente de Formação Geral As análises dos resultados de desempenho dos estudantes de Tecnologia em Gestão Pública nas duas questões discursivas relativas à Formação Geral encontram -se na Tabela 3.13 e no Gráfico 3.6.
Na Tabela 3.13, observa -se que a nota Média nesse conjunto de questões foi próxima da obtida nas objetivas.
Os estudantes, de todo o Brasil, obtiveram, em Formação Geral, Média 55,0 nas questões objetivas e 51,5 nas questões discursivas.
No entanto, pode -se notar um aumento do Desvio padrão de 21,0 nas questões objetivas do Componente de Formação Geral dos alunos de todo o Brasil, para 24,1 nas questões discursivas do mesmo componente.
A maior Média foi obtida na região Norte ( 58,3 ), e a menor, na região Sul ( 49,9 ).
A Mediana de todo o Brasil, neste componente, foi 58,5, na região Sul ( 56,5 ) a Mediana foi menor.
Nas demais regiões foi maior : Norte ( 62,0 ) ; Nordeste ( 63,0 ) ; Sudeste ( 62,3 ) ; e Centro-Oeste ( 60,0 ).
A nota Máxima ( 96,5 ) foi obtida na região Sudeste, enquanto a menor nota Máxima, foi obtida na região Norte ( 87,5 ).
A nota Mínima ( 0,0 ) foi a mesma em todas as regiões do Brasil, sem exceção.
O Gráfico 3.6 representa a distribuição das notas nas questões discursivas do Componente de Formação Geral.
A moda desta distribuição ocorre no intervalo ( 60 ; 70 ].
Destaca -se, também, o intervalo [ 0 ; 10 ] com distribuição acima de 10 \% do total de notas e representa um máximo local, sendo que no intervalo [ 0 ; 10 ] se inclui, além da nota zero, a frequência de alunos que deixaram este tipo de questão em branco.
Na sequência, os resultados verificados para cada uma das questões discursivas de Formação Geral serão apresentados, estabelecendo -se relações com os conteúdos abordados em cada uma delas.
Os comentários da Banca de docentes corretores, a respeito do observado na correção das respostas dos estudantes, suas impressões e conclusões serão apresentados junto à análise de cada questão.
Cumpre esclarecer que, tendo em vista que as questões discursivas de Formação Geral são padronizadas, ou seja, constam de todas as provas, os comentários da Banca são os mesmos para todas as carreiras acadêmicas, sendo direcionados a todos os estudantes que participaram do Enade/2015.
A seguir, serão analisados os desempenhos da Área de Tecnologia em Gestão Pública nas duas questões discursivas de Formação Geral do Enade/2015, comparando -se os resultados obtidos com comentários para cada questão. 3.3.1.1 Análise de Conteúdo da Questão Discursiva 1 do Componente de Formação Geral Os dados de Tecnologia em Gestão Pública, obtidos a partir das respostas à questão 1, encontram -se na Tabela 3.14 e no Gráfico 3.7.
Nessa questão – de melhor desempenho dentre as duas de Formação Geral – os alunos, de todo o Brasil, tiveram Média 56,5.
A maior Média para a questão 1 foi obtida na região Norte ( 63,8 ), e a menor, na região Sul ( 55,1 ).
Quanto à variabilidade das notas, o Desvio padrão de todo o Brasil foi 26,4.
O menor Desvio padrão foi obtido na região Norte ( 22,6 ), e o maior Desvio padrão foi obtido na região Centro-Oeste ( 28,3 ).
A Mediana ( 65,0 ) e a nota Mínima ( 0,0 ) foram as mesmas para todas regiões do Brasil, exceto a Mediana na região Norte ( 70,0 ).
A nota Máxima ( 100,0 ) foi obtida nas regiões Sul e Centro-Oeste, enquanto nas demais regiões foi 95,0.
O Gráfico 3.7 mostra a distribuição das notas na questão discursiva 1 do Componente de Formação Geral.
Observa -se que a maior frequência corresponde aos alunos que obtiveram nota dentro do intervalo ( 70 ; 80 ], correspondendo à moda da distribuição.
Destacam -se, também, os alunos que deixaram a questão em branco, com quase 15 \% do total. 3.3.1.2 Comentários sobre a correção de Conteúdo das respostas à Questão Discursiva 1 O enunciado da Discursiva 1 de Formação Geral estava claro e indicava um tema de extrema relevância, qual seja : a busca pela educação universal, especificamente, que apontava um caso onde ocorreu a discriminação de acesso ao ensino formal atingindo uma menina, representando a categoria das mulheres discriminadas em tantos campos da vida e, mais ainda, em sociedades que cultuam o poder soberano dos homens.
A linguagem utilizada foi clara, indicando a existência dos fatos vinculados à luta da menina Malala pelo acesso ao ensino formal das jovens do seu país, e o confronto ideológico que gerou o ataque à sua integridade física.
Diante da repercussão internacional, houve o reconhecimento pelo prêmio Nobel, e a ONU resolveu apoiá -la como símbolo da luta pela igualdade entre os sexos buscando a universalização da educação.
O conteúdo está integrado às Diretrizes, vez que o acesso ao ensino formal está vinculado ao direito constitucional à educação, que, também, é protegido pelo Estatuto da Criança e do Adolescente e nos pactos de Direitos Humanos aderidos pelo Brasil.
Por outro lado, o texto demandava uma reflexão sobre a realidade dos indivíduos que não têm acesso à educação e a situação das mulheres excluídas desse processo por segmentos da sociedade.
Como se tratava de texto contendo uma ampla gama de possibilidades, era natural a ocorrência de interpretações diversas das esperadas.
Releva notar que o enunciado pretendeu, primordialmente, identificar a questão de gênero, apontando para as dificuldades que as mulheres possuem, não somente no Brasil, mas em vários países do mundo.
Ao analisar as provas realizadas, constatamos que no item ‘ a ’ poucos foram os que identificaram os estatutos normativos indicados no padrão de resposta.
Essa ausência pode ser justificada tendo em vista que, no enunciado indicado, não existia elemento que pudesse ser utilizado pelo concluinte para embasar sua resposta nesse caminho.
As respostas foram mais gerais, identificando que a educação constitui um item fundamental para o desenvolvimento do indivíduo e que as mulheres sofrem mais com as limitações impostas por países com regime patriarcal, que discriminam o sexo feminino e, com frequência, proíbem as mulheres de frequentar a escola.
Deve -se ter em vista que a indicação de resposta que vinculava o texto aos estatutos normativos representou um raciocínio próprio dos estudantes de Direito que, naturalmente, encaminharam suas respostas para a fundamentação nessa linha identificando, não somente a Constituição, mas também leis, regulamentos, etc.
Tal linha de argumentação não foi objeto da formação dos estudantes das demais áreas.
A maior parte dos concluintes estabeleceu o tema educação formal associado à busca pelo conhecimento, como acesso ao ensino, indicando a independência intelectual e o desenvolvimento da consciência crítica, limitando as manipulações e permitindo amplo acesso ao mercado de trabalho, traduzindo uma libertação econômica e ascensão social.
No tocante ao item ‘ b ’, restou claro que a maioria esmagadora das respostas confluiu para as indicações do padrão apresentado, sendo relevante o tema de igualdade de gênero, além de candente para as concluintes do sexo feminino que desenvolveram os textos de forma concatenada, fundamentada e, muitas vezes, lamentando que, nos dias atuais, ainda se discuta a necessidade de igualar os gêneros na sociedade.
Nessa linha, foram identificadas respostas, consoante com o padrão apresentado, especialmente nos seguintes temas :  violência física e psicológica contra a mulher, incluindo a Lei Maria da Penha, no caso específico do Brasil ;  tolerância/intolerância a vestimentas, trajes, comportamentos socialmente estereotipados ;  aspectos socioculturais que impõem à mulher uma condição de submissão na sociedade, tais como a impossibilidade de a mulher manifestar seus desejos e posicionamentos em algumas culturas, entre outros ;  igualdade/desigualdade de gênero, por exemplo no mercado de trabalho, em relação à desigualdade salarial ;  ideais de Liberdade, Igualdade e \textbf{Fraternidade} : referência a esses ideais como possibilidade de equilibrar as relações de poder entre homens e mulheres.
Em muitos casos, também, os concluintes fizeram um liame entre o acesso à educação e a melhora das relações entre homens e mulheres ao permitir que a conclusão de maior grau de educação formal permita o acesso a cargos públicos relevantes, empregos bem remunerados e posições de poder, tanto na área pública quanto na área privada.
Como as opções do segundo item foram mais amplas, o padrão de resposta identificou as linhas de pensamento adotadas pelos concluintes.
A maior parte das respostas identificou, claramente, o problema da relação de gênero e suas consequências na sociedade.
Também foi frequente constatarem que existe efetiva melhora na situação das mulheres no Brasil, mas que tal movimento não ocorre em todo o mundo.
Aliás, o exemplo de Malala conduziu a uma reflexão em âmbito mundial, e os concluintes apresentaram suas análises críticas à situação da mulher em países que não admitem sequer a sua inserção no plano educacional e, muito menos, no mercado de trabalho.
A maior parte dos concluintes apresentou concatenação lógica nas suas respostas e uma minoria, pouco expressiva, não demonstrou esse domínio de 3.3.1.3 Análise de Conteúdo da Questão Discursiva 2 do Componente de Formação Geral A Tabela 3.15 mostra que o desempenho médio dos estudantes na questão discursiva 2 ( média 49,0 ) foi inferior ao obtido na questão discursiva 1 ( média 56,5 ).
A região Norte foi aquela cuja Média, nessa questão, foi maior ( 56,4 ), e a de menor Média foi a região Sul ( 47,3 ).
Quanto à variabilidade das notas, o Desvio padrão de todo o Brasil foi 26,6, superior ao obtido na questão discursiva 1 ( 26,4 ).
O maior desvio nessa questão foi obtido na região Centro-Oeste ( 28,5 ), enquanto o menor foi obtido na região Norte ( 24,5 ).
A Mediana de todo o Brasil foi 59,0, a mesma da região Centro-Oeste, enquanto nas demais regiões, a Mediana foi : 61,0 nas regiões Norte e Nordeste ; 60,0 na região Sudeste ; e 58,0 na região Sul.
A nota Máxima ( 100,0 ) foi obtida na região Sudeste, enquanto a menor Máxima, foi obtida na região Centro-Oeste ( 97,0 ).
As notas mínimas ( 0,0 ) foram as mesmas em todas as regiões do país, sem exceção.
O Gráfico 3.8 mostra a distribuição das notas na questão discursiva 2 do Componente de Formação Geral.
Observa -se que a maior frequência corresponde aos alunos que obtiveram nota situada no intervalo ( 60 ; 70 ], correspondendo à moda da distribuição.
Destacam -se, também, os alunos que deixaram a questão em branco, um máximo local com 15 \% do total de alunos.
Nota -se, ainda, que as notas ficaram um pouco mais distribuídas ao longo dos intervalos em comparação à questão discursiva de número 1. 3.3.1.4 Comentários sobre a correção de Conteúdo das respostas à Questão Discursiva 2 O enunciado da Discursiva 2 estava, igualmente, claro, indicando um tema de extrema relevância, qual seja : o reconhecimento de manifestações artísticas nascidas na periferia, como integrante da cultura reconhecida pela sociedade.
A linguagem utilizada foi clara, indicando a análise dos fatos e suas consequências no cotidiano dos partícipes e o confronto com outras visões culturais que são consideradas mais adequadas pela sociedade civil.
O conteúdo está integrado às Diretrizes, vez que o aspecto cultural e sua inter-relação com o conhecimento haurido no nível superior é essencial para a compreensão de mundo do universitário concluinte, com reflexo na sua atuação profissional.
O texto demanda uma reflexão sobre a realidade dos indivíduos que não têm acesso aos itens básicos da sociedade, como educação, saúde, cultura tradicional e segurança e que, apesar disso, criam alternativas para o seu lazer.
A incorporação do funk como manifestação cultural e artística da sociedade brasileira encontrou, basicamente, duas correntes de resposta :  os que aquiescem positivamente à integração ;  os que rejeitam, acerbamente, tal possibilidade.
No meio, uma corrente que somente admite a integração com modificações na conduta dos integrantes do movimento, com a exclusão de letras exclusivamente com apelos sexuais e incitação à violência.
A interpretação positiva indica a origem da música na periferia, notadamente nas comunidades cariocas e sua projeção para o Brasil e o mundo.
Nessa perspectiva, surge o reconhecimento como cultura popular comparada ao samba, axé e outros ritmos também nascidos no Brasil.
Também ocorre comparação com o movimento hip-hop americano.
Já a análise negativa prioriza a vinculação do ritmo com traficantes, prostitutas, vilipêndio às mulheres, falta de controle quanto ao som, local dos eventos, transações sexuais em público, uso frequente de drogas, violência \textbf{exacerbada}.
Faltam, para essa corrente, elementos básicos da cultura, tratando -se de música sem autonomia como movimento de mudança ou aspecto próprio.
A corrente intermediária reconhece o vínculo originário com as comunidades dominadas por quadrilhas de traficantes, mas admite o seu reconhecimento de acordo com a mudança dos atores que transitam na modalidade, como músicos, \textbf{intérpretes}, para adequar as letras das canções expurgando os aspectos considerados \textbf{nocivos} como o apelo sexual, as indicações positivas para a violência e apologia às drogas.
Um aspecto bastante criticado nas músicas dessa modalidade foi a exploração da mulher, entendida como um objeto de satisfação masculina.
É relevante tal constatação, pois a questão pertinente à Malala reflete a luta pela valorização da mulher.
Os concluintes compreenderam a ideia que foi problematizada e escolheram uma estratégia adequada para solucioná -la.
O caráter dinâmico da cultura foi comumente expressado, não através do enfrentamento abordado no texto, mas mediante a indicação ( de forma rasa e genérica ) das dificuldades e consequente aceitação social de manifestações culturais, antes marginalizadas e que hoje integram o acervo cultural nacional, tais como : samba, grafite, capoeira e forró.
Quanto às reflexões de Laraia, a grande maioria reconheceu o funk como manifestação cultural, porém, o identificou mais na sua expressão musical, no balanço de seu ritmo, tecendo contundentes críticas, e mesmo forte discordância quanto ao conteúdo e mensagem das letras quando remetem à depreciação da imagem feminina ou induzem à apologia ao crime.
De resto, os textos abordaram bem mais o item dois do padrão de Respostas.
Uma minoria enfrentou os temas mudança e dinamismo cultural.
As respostas da corrente positiva indicam que o preconceito sempre acompanhou as manifestações populares como o samba, na sua origem.
Admitem que diante da miscigenação surgem movimentos culturais de várias naturezas, o que inclui o funk.
Também indicam a dificuldade de lazer dos membros das comunidades o que faz surgirem eventos para preencher essa necessidade.
Nesse âmbito, apontam a música como alternativa de ascensão social e integração aos bens de consumo.
As respostas da corrente negativa vinculam a música ao tráfico de drogas cujos elementos comandam as comunidades, propiciando que os bailes sejam mecanismos de venda e consumo dessa mercadoria ilícita.
Na mesma toada, indicam a vinculação com a exploração sexual, principalmente das mulheres maiores e menores de idade.
Indicam aspectos imorais nas letras.
Apresentam os locais dos bailes como inadequados e provocadores da \textbf{desordem} pública.
Recusam que músicas com discriminação às mulheres sejam consideradas como cultura.
A corrente intermediária indica a possibilidade de adoção do funk como cultura desde que ocorram modificações nas letras e no comportamento dos atores ( músicos, letristas e público ), retirando os aspectos negativos já apontados.
Fala -se em “ funk do bem ”.
A maior parte das respostas apresentadas indica que os concluintes têm conhecimento dos conflitos sociais e, no caso da questão do funk, conseguem identificar a origem popular da música.
A partir daí é que ocorre a cisão entre os que defendem o movimento como cultura, aduzindo a ascensão social dos artistas envolvidos dos que vinculam a música com atos ilícitos, rejeitando sua inserção no âmbito cultural.
Em algumas provas, há dificuldade de concatenação lógica de raciocínio, mas a maior parte consegue desenvolver o texto de forma razoável.
Essa questão permite visualizar a divisão social com os pensamentos \textbf{liberais} e conservadores, sendo bastante claros nas respostas dos concluintes. 3.3.1.5 Análise de Língua Portuguesa das Questões Discursivas do Componente de Formação Geral Os dados de Tecnologia em Gestão Pública, obtidos a partir das respostas às questões discursivas do Componente de Formação Geral, no que tange à Língua Portuguesa, encontram -se na Tabela 3.16 e no Gráfico 3.9.
Nesse aspecto, os alunos, de todo o Brasil, tiveram Média 50,9.
A maior Média com respeito à Língua Portuguesa foi obtida na região Norte ( 57,7 ), e a menor, na região Sul ( 49,3 ).
Quanto à variabilidade das notas, o Desvio padrão de todo o Brasil foi 24,5.
O menor Desvio padrão foi obtido na região Norte ( 21,1 ) e o maior Desvio padrão foi obtido na região Centro-Oeste ( 25,7 ).
A Mediana das notas de Língua Portuguesa foi 57,5 para o brasil como um todo, sendo a maior e a menor Mediana obtidas, respectivamente, nas regiões Nordeste ( 62,0 ) e Sul ( 55,5 ).
A nota Máxima para todo o Brasil foi de 98,0, com, pelo menos, um aluno tirando essa nota nas regiões Nordeste e Sul ( 98,0 ), enquanto a menor nota Máxima foi obtida na região Norte ( 88,5 ).
Além disso, a nota Mínima foi zero em todas as regiões do país, sem exceção.
O Gráfico 3.9 mostra a distribuição das notas de Língua Portuguesa do Componente de Formação Geral.
Observa -se que a maior frequência ( quase 25 \% ) corresponde aos alunos que obtiveram nota no intervalo ( 60 ; 70 ].
Destacam -se, também, os alunos que deixaram a questão em branco, um máximo local com quase 12,5 \% do total de alunos 3.3.1.6 Comentários sobre a correção das respostas de Formação Geral com respeito à Língua Portuguesa Ao encaminhar as questões 1 e 2 na direção da produção de um texto dissertativo, espera -se que o participante utilize seus conhecimentos sobre o assunto e estruture seus textos de acordo com as exigências do registro formal próprio dessa situação comunicativa.
Essa configuração determina exigências quanto : à adequação da seleção vocabular, ao desenvolvimento do conteúdo, à estruturação sintática dos períodos, à organização lógica das ideias, à utilização de procedimentos de encadeamento textual e referenciação, à obediência às exigências morfossintáticas próprias da modalidade escrita da norma-padrão, ao respeito às regras ortográficas e às regras de acentuação gráfica.
O padrão de resposta utilizado na avaliação das questões 1 e 2 considerou os aspectos relevantes ao bom desempenho linguístico como competências distintas, de modo a permitir um mapeamento detalhado do domínio dos recursos disponíveis na Língua Portuguesa para a comunicação escrita formal.
Com base nesse objetivo, foram avaliados os seguintes aspectos : a ) Estruturação textual condizente com o gênero solicitado ( texto dissertativo ) e o modo de organização textual expositivo adequado ao gênero – essa competência envolve :  distribuição do conteúdo do texto em parágrafos, de modo a garantir a sua organização temática ;  estruturação sintática condizente com o padrão da modalidade escrita formal da língua portuguesa de modo a garantir a clareza necessária ;  utilização de operadores discursivos que contribuam para a progressão temática do texto, estabelecendo relações lógicas entre as ideias apresentadas, tanto do ponto de vista intrafrasal, como do interfrasal ;  utilização de procedimentos de referenciação lexical e pronominal que permitam a retomada de referentes textuais ;  utilização de sinais de pontuação que contribuam para a organização lógica da frase e do texto.
Espera -se, portanto, que o participante recorra a procedimentos linguístico-discursivos para organizar seu texto, permitindo o encadeamento lógico entre suas partes de forma a garantir a progressão e a coerência textual.
Isso significa que os seguintes procedimentos foram penalizados, de acordo com o padrão de resposta proposto :  elaboração de frases fragmentadas que comprometam a estrutura lógico-gramatical do texto ;  sequência justaposta de ideias sem encaixamentos sintáticos, reproduzindo hábitos da oralidade ;  elaboração de frase com apenas oração subordinada, sem oração principal ;  emprego equivocado do conector ( preposição, conjunção, pronome relativo, alguns advérbios e locuções adverbiais ) que não expresse a relação lógica adequada entre dois trechos do texto e prejudique a compreensão da mensagem ;  repetição ou substituição inadequada de palavras sem se valer dos recursos oferecidos pela língua ( pronome, advérbio, artigo, sinônimo ) ;  utilização inadequada dos sinais de pontuação comprometendo a clareza textual. b ) Respeito às convenções ortográficas da norma-padrão da Língua Portuguesa – essa competência envolve o domínio das regras de acentuação gráfica e da grafia padrão das palavras ( com ausência de abreviaturas próprias da linguagem da internet ), de acordo com as convenções estabelecidas pela legislação em em vigor e consubstanciadas no Vocabulário Ortográfico da Língua Portuguesa, editado pela Academia Brasileira de Letras ( com aceitação da legislação anterior, no caso das regras relativas ao uso do hífen e da acentuação gráfica, já que o exame foi realizado no mês de \textbf{novembro} de 2015 ).
A competência relativa ao respeito às convenções ortográficas envolve : • grafia correta das palavras ; • respeito às regras de acentuação gráfica ; • emprego de maiúsculas em início de frase, em nomes próprios de pessoas, lugares ou instituições ; • ausência de abreviações como p/, vc, tb, pq, tá, né, usadas muitas vezes em escrita informal e na internet ; • obediência às regras de separação de sílabas no final da linha. c ) Domínio dos diferentes aspectos morfossintáticos próprios da modalidade escrita formal da norma-padrão da Língua Portuguesa – os itens relacionados a essa competência são : a concordância nominal, a concordância verbal, a regência nominal, a regência verbal, a flexão nominal, a flexão verbal, a correlação entre os tempos verbais, a colocação pronominal.
O domínio dessa competência envolve : • flexão do verbo para estabelecer concordância de número com o sujeito da frase ; • flexão do artigo, do adjetivo e do pronome para concordar em número e em gênero com o substantivo a que se referem ; • respeito à regência nominal e à verbal, utilizando a preposição adequada depois de um substantivo, de um verbo ou de um adjetivo ; emprego da preposição antecedendo o pronome relativo e emprego adequadamente os pronomes relativos “ cujo(a) ” e “ onde ” ; • emprego adequado do acento grave indicador da crase entre uma preposição e um artigo ( a+a ) ; • atendimento às regras de colocação pronominal ( próclise e ênclise ), distintas dos hábitos da oralidade ou da escrita informal. d ) Seleção vocabular adequada à modalidade escrita formal da Língua Portuguesa, exigida pela situação comunicativa – essa competência envolve a precisão na utilização do vocabulário relacionado à temática solicitada pela questão ; a ausência de marcas da oralidade, como termos de sentido muito genérico ( “ coisa ”, “ negócio ”, “ você ” ) e termos de registros mais informais ( como gírias, jargões, frases feitas, ditados populares, termos regionais ).
Assim, espera -se que o participante respeite a adequação vocabular não empregando gírias ou expressões coloquiais, evite repetição desnecessária de palavras e utilize um vocabulário mais formal, como solicitado por um texto dissertativo.
A escolha dessas competências para subsidiar o processo de avaliação apoia -se na concepção de que, no desempenho dos graduandos, a modalidade escrita tem apresentado uma intensa simplificação, originada no padrão da modalidade oral da Língua Portuguesa.
No caso do texto de base dissertativa, inscrito em um registro formal, a distância entre as duas modalidades é ainda maior, o que provoca situações de hipercorreção ( desvios provocados pela utilização inadequada de uma regra da norma-padrão ) e de truncamentos sintáticos ( estruturas frasais incompreensíveis devido à complexidade sintática própria da modalidade escrita ).
Observam -se, então, os seguintes aspectos que marcam essa distinção entre as duas modalidades, devido à excessiva simplificação da modalidade falada : a ) redução drástica de estruturas subordinadas, compensada pelo aumento na frequência de estruturas coordenadas e absolutas, por um lado, ou pela elaboração de estruturas truncadas pelo excesso de ideias sem a devida conexão subordinativa ; b ) redução no uso de operadores argumentativas para expressar relações lógicas essenciais à construção do texto, substituídas pela exigência de inferência por parte do interlocutor para suprir a sua ausência ; c ) redução do uso do subjuntivo, ao lado da ampliação do uso do indicativo combinado a estruturas frasais coordenadas ou absolutas ; d ) empobrecimento do processo de referenciação, com a repetição de pronomes ou nomes ; e ) simplificação extrema da marcação da categoria tempo na morfologia verbal ; f ) falta de domínio de vocabulário mais abstrato e de maior complexidade, essencial ao desenvolvimento do processo dissertativo ; g ) redução drástica no emprego da acentuação gráfica, processo intensificado pela intensa utilização das redes sociais.
Os aspectos macroestruturais da elaboração do texto não foram avaliados neste processo porque dizem respeito à avaliação do conteúdo.
Portanto, a banca de Formação Geral, composta por profissionais de diferentes áreas do conhecimento, se encarregou da avaliação do atendimento ao solicitado no enunciado das questões do ponto de vista do desenvolvimento do conteúdo.
A avaliação do desempenho linguístico considerou, portanto, três grandes grupos de competências, segundo os aspectos explicitados anteriormente : a ) domínio das convenções ortográficas : grafia de vogais e consoantes, uso de maiúsculas e minúsculas, emprego do hífen e acentuação gráfica ; b ) domínio dos procedimentos de estruturação textual do ponto de vista microestrutural : organização interna dos períodos, emprego de conectores para a articulação lógica entre os períodos e entre os parágrafos, emprego de marcas de referenciação lexical e pronominal ; utilização dos sinais de pontuação que contribuem para a organização lógica da frase. c ) domínio das regras de caráter morfossintático estabelecidas como modelares do ponto de vista da modalidade escrita formal da norma-padrão da Língua Portuguesa : concordância nominal e verbal, regência nominal e verbal, colocação pronominal, flexão nominal e verbal, correlação entre tempos e modos verbais, ausência de marcas de oralidade.
A seleção vocabular adequada à modalidade escrita formal da Língua Portuguesa foi incorporada à essa última competência, tendo em vista a intersecção entre as duas do ponto de vista das exigências do registro formal da modalidade escrita da norma-padrão.
Apresenta -se, a seguir, o padrão de resposta aprovado pelo Inep e já utilizado na avaliação do desempenho linguístico das questões 1 e 2 de Formação Geral no Enade/2013 e no Enade/2014.
Às competências, reunidas nos três grupos descritos anteriormente, foram atribuídos os seguintes \textbf{pesos} relativos : aspectos ortográficos ( 20 \% ) ; aspectos textuais ( 40 \% ), aspectos morfossintáticos e vocabulares ( 40 \% ).
Com base na avaliação do processo desenvolvido nos dois anos anteriores, utilizaram -se quatro níveis para a avaliação do desempenho linguístico nas três competências.
O nível zero ficou reservado para casos especiais, como fuga total ao tema, palavras soltas, produção de uma frase incompleta, entre outros.
A avaliação das respostas elaboradas pelos participantes revelou um resultado coerente na comparação entre as duas questões.
Entretanto, em relação aos resultados do Enade/2014, observou -se uma melhora do domínio de todas as competências, talvez devido ao perfil dos cursos analisados, como Direito e Jornalismo, que pressupõem a utilização profissional da linguagem.
Observou -se uma diferença no desempenho dos participantes em relação às duas questões discursivas, provavelmente em função dos temas solicitados : - a Questão 1 solicitava que o participante desenvolvesse um texto dissertativo sobre o significado da premiação de Malala, abordando dois aspectos : o direito das jovens à educação formal e as relações de poder entre homens e mulheres no mundo.
Esse encaminhamento favoreceu o desenvolvimento mais articulado do texto, o que se refletiu em melhor desempenho dos aspectos vocabulares e textuais. - a Questão 2 solicitava que o participante desenvolvesse um texto dissertativo para discutir a questão do reconhecimento do funk como legítima manifestação artística e cultural do povo brasileiro.
Por ser um tema muito polêmico, o tema do preconceito sócio-cultural-econômico favoreceu a construção de textos mais críticos, dividindo o conjunto de participantes em dois grupos : os que defendiam esse tipo de manifestação artística como legítima e os que a criticavam de maneira veemente, chegando, inclusive, a utilizar muitos termos de baixo calão.
Essa última atitude teve como consequência a elaboração de textos mal desenvolvidos do ponto de vista sintático-discursivo, com muitos truncamentos.
Os resultados revelaram essa diferença.
Ao analisar os dados de cada questão, constata -se que a Questão 1 teve maior número de questões de médio ou bom desempenho, em relação à Questão 2.
Por outro lado, a Questão 2 apresentou um índice maior de notas baixas em relação à Questão 1.
Os resultados da avaliação correspondem aos seguintes aspectos observados em cada competência : a ) aspectos ortográficos : o desempenho dos participantes revelou uma diferença muito grande nos dois aspectos analisados : baixo índice de desvios ortográficos e grande índice de desvios de acentuação.
Em vários casos, ocorre ausência completa de acentuação gráfica.
Os resultados revelam que a tendência dominante entre os universitários brasileiros é a eliminação da acentuação gráfica, talvez motivada pelos hábitos relacionados às redes sociais e pela ausência de esclarecimento dos meios de comunicação, das autoridades e das escolas sobre as decisões do Acordo Ortográfico da Língua Portuguesa de 1990.
Os casos mais sistemáticos de eliminação do acento indicador da sílaba tônica são :  palavras proparoxítonas ( “ proximos ”, “ politicos ” ) ;  palavras paroxítonas terminadas em ditongo crescente ( “ necessario ”, “ noticiarios ”, “ individuo ”, “ dependencia ”, “ varios ”, “ propria ” ) ;  palavras oxítonas ( “ ninguem ”, “ esta ”, “ ate ”, “ tambem ”, “ alvara ” ) ;  uso indevido do acento gráfico em palavras como “ genêro ”, “ melâncolia ”, por exemplo.
Quanto ao domínio das convenções relativas à grafia das palavras, observam -se desvios como : I. hipercorreção pela escolha de “ e ” no lugar de “ i ” ou de “ i ” no lugar de “ e ”, por influência de hábitos da oralidade ( “ descriminação ” por “ discriminação ”, “ entervenção ” por “ intervenção ”, “ esteriotipada ” por “ estereotipada ” ) ; II. eliminação do “ r ” marcador do infinitivo verbal ( “ esta ” no lugar de “ estar ”, “ estuda ” no lugar de “ estudar ” ).
III. eliminação de uma sílaba na palavra “ educação ”, que foi grafada “ edução ” por um grande número de participantes.
IV. outros casos de desvios de grafia relacionados à variação diastrática podem ser observados em “ estrupo ”, “ subjulgadas ”, “ precoseito ”, “ soubre ”, “ vecendo ”, “ indesencia ”, “ apolojia ”, “ fulga ”.
V. desvios de segmentação podem ser atribuídos à confusão entre o \textbf{morfema} gramatical “ mos ” e o pronome oblíquo “ nos ”, em “ chegar mos ”, “ enfrentar -mos ”.
VI. inadequação no uso da maiúscula para destacar determinadas palavras-chave do texto, como “ Violência ”, “ Brasileiros ”.
Destaque -se, também, o grande número de participantes que grafam os textos inteiramente em caixa alta.
Vale observar, também, que, ao contrário do que se esperava, não apareceram abreviaturas próprias do “ internetês ”, relacionado ao uso de redes sociais e emails. b ) aspectos textuais : esta competência é a que se revela como a mais problemática entre os participantes, porque são muitos os problemas observados, desvios acumulados durante toda a formação do estudante : sequência justaposta de ideias sem encaixamentos sintáticos ; redução drástica de estruturas subordinadas, ao lado do aumento na frequência de estruturas coordenadas e absolutas ; redução no uso de conectores para expressar relações lógicas essenciais à construção do texto, substituídas pela exigência de inferência por parte do interlocutor para suprir a sua ausência ; emprego equivocado de operadores que não estabelecem relações lógicas coerentes entre ideias do texto ; emprego inadequado do pronome relativo ( com omissão da preposição ou a utilização de pronome inadequado, como “ onde ” ) ; repetição exaustiva de termos sem a utilização de procedimentos mais sofisticados de substituição ( hiperonímias, hiponímias, nominalizações, expressões metafóricas ) ; frases fragmentadas que comprometem a estrutura lógico-gramatical ; frases formadas apenas por oração subordinada, sem oração principal. ( hiperonímias, hiponímias, nominalizações, expressões metafóricas ) ; frases fragmentadas que comprometem a estrutura lógico-gramatical ; frases formadas apenas por oração subordinada, sem oração principal.
Um importante aspecto a destacar é o baixíssimo desempenho de uma parte dos participantes em relação à estrutura formal do texto produzido, o que é extremamente preocupante ao se levar em conta que são graduandos em fase final de formação.
São frequentes os casos de desvios de estruturação frasal, com uso inadequado ou ausência de conectivos entre parágrafos e entre frases.
Em uma parte dos textos, falta um mínimo de textualidade e de domínio do registro padrão da língua.
Na verdade, observam -se relações linguísticas quase agramaticais, como as estabelecidas pela sequência de gerúndios sem o apoio de um ponto de partida para a organização das informações gramaticais e semânticas.
Observou -se que uma grande parte dos participantes não organizou as ideias em parágrafos, talvez devido ao pequeno número de linhas disponibilizadas para a resposta da questão ou, quem sabe, pela suposição de que não seria necessária essa divisão por não se tratar de um texto no modelo de uma redação dissertativo-argumentativa, como solicitado nos vestibulares.
Quanto à utilização dos mecanismos de referenciação, deve -se destacar a ocorrência, em uma boa parte dos textos, de repetições de palavras ou expressões sem a utilização de termos sinônimos ou pronomes, como seria adequado.
Outro aspecto relevante a ser destacado é a quase total ausência de operadores argumentativos, tanto intrafrasais como interfrasais, repercutindo uma tendência atual da mídia escrita.
Quanto à utilização dos sinais de pontuação, observou -se uma grande precariedade nos textos analisados.
É muito frequente a ocorrência de parágrafos sem marca interna de pontuação para separar os períodos.
Vale observar que não foi penalizada a ausência de vírgula para destacar locuções ou adjuntos adverbiais de pequena extensão deslocados de posição na frase, por ser um uso opcional.
São os seguintes os tipos de problemas encontrados : I. vírgula : utilização de vírgula para separar o sujeito e o predicado ; ocorrência de apenas uma das vírgulas para separar uma palavra, uma expressão ou uma oração encaixada ; uso de vírgula no lugar do ponto para separar ideias que constituem períodos distintos ; II. ponto e vírgula : utilização do ponto e vírgula no lugar de vírgula ; III. ponto final : ausência de ponto final para separar períodos. c ) aspectos morfossintáticos e vocabulares : os resultados são muito transparentes em relação aos aspectos mais problemáticos no desempenho dos participantes.
O desvio mais frequente, em relação à regência, é a falta do sinal indicativo da crase – isso revela que o usuário não tem consciência de que, sob a forma do termo “ a ”, existe a presença de uma preposição “ a ”, exigida pela regência do termo anterior, e de um artigo definido.
Embora em outros exames, como o Enem, a falta de crase seja penalizada em acentuação, nesta avaliação esse desvio foi considerado no âmbito dos aspectos morfossintáticos.
Outro problema relacionado à regência verbal e à nominal, encontrado frequentemente nas questões, foi a ausência de preposição antes de pronome relativo, processo generalizado na modalidade oral da língua, em situações de registro informal.
Apesar da possibilidade de que essa alteração de regência se generalize no padrão escrito da Língua Portuguesa, como já está ocorrendo até em textos jornalísticos, o não emprego da preposição foi penalizado neste processo de avaliação.
Outro desvio muito frequente diz respeito aos processos de concordância verbal e de concordância nominal.
Quanto à concordância de número, observou -se ausência de marca com sujeito posposto ou uso indevido ( uso inadequado da marca de plural comandado pelo núcleo plural da locução adjetiva, apesar de o substantivo que funciona como núcleo do sintagma nominal estar no singular ).
Uma ocorrência generalizada foi a ausência de acento circunflexo na forma plural do presente do indicativo do verbo “ ter ”, que foi considerada como um desvio na concordância verbal e não na acentuação gráfica.
Quanto à concordância de gênero, vários casos foram observados, normalmente no âmbito de sintagmas nominais longos, em que o adjetivo está afastado do substantivo.
Quanto à questão da colocação pronominal, foram poucos os desvios observados.
Apesar de serem aspectos relacionados à oralidade, concluiu -se que, no registro escrito formal, a maioria dos participantes já incorporou regras como a não introdução da frase por um pronome oblíquo e a próclise na presença de um termo atrator.
Não se adotou, entretanto, o padrão excessivamente formal descrito pelas gramáticas normativas em relação à posição do pronome oblíquo em locuções verbais, já que esse uso está muito distante da prática cotidiana, até em textos mais formais.
Quanto aos aspectos vocabulares, alguns tipos de inadequação foram observados, principalmente na Questão 2, relacionada ao funk : expressões da oralidade, seleção vocabular incompatível com o contexto, gerando situações de falta de inteligibilidade ; falta de domínio de vocabulário mais abstrato e de maior de complexidade, essencial ao desenvolvimento do texto de base dissertativa.
Ainda em função da temática desenvolvida na Questão 2, foi significativa a presença de palavras de baixa calão, por retratar um universo cultural específico.
Os textos que continham esse tipo de vocabulário foram considerados inadequados em função do registro formal exigido e foram avaliados com grau zero.
Várias marcas de oralidade foram identificadas, embora não em alta frequência : o uso do pronome relativo “ onde ” como relativo universal, falta de artigo definido antes de substantivo, repetição de palavras por falta de vocabulário, reduções como “ tá ” e “ pra ”, expressões informais. 3.3.2 Componente de Conhecimento Específico Na parte da prova relativa às questões discursivas do Componente de Conhecimento Específico ( Tabela 3.17 ), observa -se que a Média foi mais baixa do que para as questões discursivas do Componente de Formação Geral.
Enquanto no Componente de Formação Geral a Média para estudantes de Tecnologia em Gestão Pública de todo o Brasil foi 51,5, na parte de Conhecimento Específico a Média foi 39,9.
A maior Média deste componente foi obtida pelos estudantes da região Nordeste ( 47,2 ), e a menor, pelos da região Sul ( 38,0 ).
Quanto à variabilidade das notas, o Desvio padrão de todo o Brasil foi 26,3.
O maior Desvio padrão foi encontrado nas regiões Sudeste e Centro-Oeste ( 26,5 cada ), e o menor, na região Norte ( 23,1 ).
A nota Máxima ( 100,0 ) e a nota Mínima ( 0,0 ) foram obtidas por pelo menos um aluno em quase todas as regiões do Brasil, exceto a nota Máxima nas regiões Norte ( 95,0 ) e Centro-Oeste ( 93,3 ).
A Mediana do Brasil como um todo foi 41,7.
Nas regiões Norte e Sudeste a Mediana foi 48,3, na região Nordeste, foi 50,0, e nas regiões Sul e Centro-Oeste, foi 40,0.
O Gráfico 3.10 representa a distribuição das notas nas questões discursivas no Componente de Conhecimento Específico.
A moda desta distribuição ocorre no primeiro intervalo, [ 0 ; 10 ].
A análise de cada uma destas questões será feita a seguir. 3.3.2.1 Análise da Questão Discursiva 3 do Componente de Conhecimento Específico Na questão 3, cujos resultados aferidos encontram -se descritos na Tabela 3.18, a Média dos estudantes de todo o Brasil foi 30,6.
O desempenho dos estudantes de todo o Brasil nesta questão foi mais baixo dentre as três questões discursivas desse componente.
A menor Média nessa questão foi obtida pelos alunos da região Centro-Oeste ( 27,8 ), enquanto a maior Média foi obtida na região Nordeste ( 33,9 ).
Quanto à variabilidade das notas, o Desvio padrão de todo o Brasil foi 28,0.
O maior Desvio padrão foi obtido na região Sudeste ( 29,1 ), enquanto o menor foi obtido na região Sul ( 27,6 ).
As medianas ( 35,0 ), as notas máximas ( 100,0 ) e as notas mínimas ( 0,0 ) foram iguais para quase todas as regiões, as exceções foram a Mediana na região Centro-Oeste ( 20,0 ) e a nota Máxima na região Norte ( 85,0 ).
O Gráfico 3.11 mostra a distribuição das notas na questão discursiva 3, do Componente de Conhecimento Específico da área de Tecnologia em Gestão Pública.
Sem levar em consideração que quase 30 \% dos estudantes deixaram a questão em branco, essa distribuição apresenta moda no intervalo ( 30 ; 40 ], com os intervalos ( 10, 20 ], ( 50, 60 ] e ( 80, 90 ] aparecendo como máximos locais. 3.3.2.2 Comentários sobre as respostas à Questão Discursiva 3 Os conhecimentos exigidos nesta questão trataram do problema das externalidades, positivas e negativas, enquanto uma das falhas de mercado.
O conteúdo é fundamental e básico em cursos de Administração e Gestão Pública.
O enunciado, apesar de bastante claro e de apresentar uma linguagem acessível e direta, não demandou profundidade de conhecimento dos conteúdos avaliados, nem excesso de conteúdo.
A temática da questão foi apropriada para o programa de Tecnologia em Gestão Pública - TGP.
De fato, demandou conhecimento específico da teoria de economia do setor público, particularmente de uma falha de mercado, e o conceito exigido é bastante apropriado para o nível de um exame em Tecnologia em Gestão Pública.
Os dados eram pertinentes e suficientes para a solução do problema.
A questão foi elaborada de forma clara, não dando azo a interpretações outras que levassem a respostas diversas ou erradas, por erro de interpretação.
Cerca de 1/3 dos estudantes deixaram a resposta a essa questão em branco.
Muitas outras foram respondidas de maneira incompleta, mas nenhuma resposta que indicasse que o conteúdo não foi entendido por aqueles que se dedicaram a respondê -la.
A análise das provas mostrou que todos os alunos que responderam à questão foram capazes de dar um exemplo de política saneadora ou mitigadora dos efeitos negativos de uma externalidade negativa.
Ou seja, eles sabiam a qual problema se referia o enunciado e entenderam a questão.
Ademais, dos que responderam à questão, apenas pouco mais de 100 em mais de 3.000 respondentes, tiraram nota"zero", o que reafirma que a maioria dos respondentes sabia à que se referia a pergunta.
Não houve ocorrência de interpretações diferentes das esperadas pelo padrão adotado.
O espaço para apresentação da solução e o tempo da prova foram adequados.
A questão pode ser considerada média, em comparação com as demais questões de Conhecimentos Específicos, porque aborda um fenômeno conhecido.
O exemplo é, de fato, corrente, e questiona sobre uma das falhas de mercados, de mais fácil entendimento.
Contudo, como o enunciado pedia a conceituação de externalidades positivas e negativas, poucos respondentes foram capazes de definir externalidades, indicando que o nível de abstração demandado para apresentar o conceito foi de maior do que se esperava.
Com respeito à diversidade das respostas encontradas, a maioria demonstrou compreensão do enunciado.
Contudo, apesar de saberem discutir o assunto e apresentar exemplos e políticas apropriadas, os estudantes demonstraram não saber conceituar o problema formalmente.
Poucos alunos foram capazes de conceituar o que são externalidades de maneira formal.
Neste sentido, as respostas foram incompletas.
Como o exemplo usado, poluição oriunda de emissão de gases combustíveis, é bastante conhecido, e todos estavam familiarizados com o tema da poluição causada pelo excesso de carros circulando nas grandes cidades, os alunos conseguiram elaborar sobre o problema com fluência.
Quase todos os alunos parafrasearam parte do enunciado, na maioria das vezes, de forma construtiva e integrada na resposta.
A linguagem foi apropriada do ponto de vista dos termos específicos da Área, mas os alunos não fizeram uso de linguagem cientifica ou sofisticada, ou seja, usaram a linguagem popular e corrente para descrever o fenômeno.
Não que isto pese no valor da resposta, uma vez que, gestão pública e suas questões, pertencem ao cotidiano de todos nós.
As notas dos estudantes foram distribuídas em três níveis : fracas ; medianas e boas.
A maioria dos alunos respondeu diretamente ao segundo ou terceiro itens da questão ( itens ‘ b ’ e ‘ c ’ ).
Apenas exemplificaram externalidades e deram exemplos de políticas públicas para corrigir o problema da externalidade negativa.
Pouquíssimos conseguiram ou tentaram conceituar externalidade.
Alguns alunos, que tentaram conceituar, confundiram o conceito com o exemplo.
Não houve muita diversidade de soluções apresentadas, e quase todas as respostas ao item ‘ c ’ ( políticas públicas saneadoras ) sugeriram o uso de transporte alternativo, particularmente, transportes coletivos, e o incentivo ao uso de bicicletas.
O principal problema foi a ausência de respostas sobre conceitos, demonstrando uma falha na formação dos alunos, ou respostas incompletas.
Nas respostas consideradas fracas, alguns não conceituaram, não deram exemplos ou não indicaram políticas públicas, fazendo comentários e usando os próprios termos do enunciado com exemplos equivocados ou errados.
Outros não obtiveram sucesso ao abordar o assunto da ótica solicitada pelo enunciado, mesmo utilizando a terminologia desse enunciado.
Entre aqueles que conseguiram obter dois pontos, estão os que citaram, pelo menos, uma política pública visando mitigar, por exemplo, o problema da poluição.
Nas respostas consideradas medianas estão aqueles que citaram apenas uma externalidade e uma solução de política pública.
Ou, ainda, estão aqueles que apontaram duas propostas de políticas públicas para mitigar externalidades negativas, mas não conceituaram ou exemplificaram.
E alguns outros apontaram exemplos apenas de externalidades negativas e propostas de políticas públicas para mitigar apenas essas externalidades.
Nas respostas consideradas boas estão aqueles que souberam apontar exemplos de externalidades positivas e negativas e apresentaram propostas de políticas para mitigar essas externalidades, mas não as conceituaram.
E entre os que obtiveram nota máxima estão os que apresentaram a questão completa, abordando o conceito, exemplificando as externalidades positivas e negativas e sugerindo políticas públicas mitigadoras dos problemas causados das externalidades negativas.
A principal deficiência na formação dos alunos em relação ao conteúdo exigido e à capacidade de expressão desses conhecimentos foi a inabilidade de conceituar o fenômeno analisado e abstraí -lo do exemplo.
Não saber definir o conceito de externalidade demonstra que o aluno não conseguiu abstrair a ideia deste.
Ou seja, só conseguiu lidar com o aspecto prático, ligado à realidade cotidiana.
Como ensina José Carlos Libâneo, citando Vygostky e Davydov, conceitos científicos deveriam estar incorporados nos processos de pensamento teórico, mas ser engendrados com a prática.
Lidar com o conteúdo só na teoria não faz sentido, e o pensamento ocorre também durante a formação e operação de conceitos.
Estes processos deveriam estar integrados entre si, formando uma unidade.
O aluno deveria ter incorporado os conceitos científicos nos processos de pensamento e vice-versa.
Mas, na maioria das respostas, parece que o aluno nunca aprendeu o conceito per se, e, sim, somente o fenômeno.
Uma recomendação para professores que ministram esta disciplina seria enfatizar o conceito de externalidade e mostrar sua pertinência e aplicação.
A correção foi adequada e levou em considerou todas as respostas apresentadas, pontuando toda contribuição feita pelos alunos.
Os critérios de distribuição de pontos eram precisos, o que facilitou a correção. 3.3.2.3 Análise da Questão Discursiva 4 do Componente de Conhecimento Específico A Tabela 3.19 contém as informações relativas à questão 4 do conjunto de questões do Componente de Conhecimento Específico.
O desempenho dos estudantes nessa questão foi superior ao das questões 3 e 5.
A Média geral do Brasil foi 48,0, sendo a menor Média registrada na região Sul ( 45,8 ), e a maior, na região Nordeste ( 58,4 ).
As medianas ( 60,0 ), as notas máximas ( 100,0 ) e as notas mínimas ( 0,0 ) foram iguais para todas as regiões, sem exceção.
O Gráfico 3.12 representa a distribuição de notas da questão discursiva 4, do Componente de Conhecimento Específico.
Essa distribuição tem moda no intervalo ( 50 ; 60 ].
Lidar com o conteúdo só na teoria não faz sentido, e o pensamento ocorre também durante a formação e operação de conceitos.
Estes processos deveriam estar integrados entre si, formando uma unidade.
O aluno deveria ter incorporado os conceitos científicos nos processos de pensamento e vice-versa.
Mas, na maioria das respostas, parece que o aluno nunca aprendeu o conceito per se, e, sim, somente o fenômeno.
Uma recomendação para professores que ministram esta disciplina seria enfatizar o conceito de externalidade e mostrar sua pertinência e aplicação.
A correção foi adequada e levou em considerou todas as respostas apresentadas, pontuando toda contribuição feita pelos alunos.
Os critérios de distribuição de pontos eram precisos, o que facilitou a correção. 3.3.2.3 Análise da Questão Discursiva 4 do Componente de Conhecimento Específico A Tabela 3.19 contém as informações relativas à questão 4 do conjunto de questões do Componente de Conhecimento Específico.
O desempenho dos estudantes nessa questão foi superior ao das questões 3 e 5.
A Média geral do Brasil foi 48,0, sendo a menor Média registrada na região Sul ( 45,8 ), e a maior, na região Nordeste ( 58,4 ).
As medianas ( 60,0 ), as notas máximas ( 100,0 ) e as notas mínimas ( 0,0 ) foram iguais para todas as regiões, sem exceção.
O Gráfico 3.12 representa a distribuição de notas da questão discursiva 4, do Componente de Conhecimento Específico.
Essa distribuição tem moda no intervalo ( 50 ; 60 ]. 3.3.2.4 Comentários sobre as respostas à Questão Discursiva 4 A pergunta apresentou um enunciado claro sobre a relação entre gestão ética e transparência.
Este conteúdo é particularmente pertinente ao estudo de Gestão Pública.
Elaborar sobre o assunto e ser capaz de relacionar as características essenciais do conceito de gestão ética e de transparência são algumas das competências e habilidades almejadas nas Diretrizes Curriculares da Área.
O conteúdo é fundamental e básico dos cursos de Administração e Gestão Pública.
A linguagem usada foi acessível e direta.
A questão era clara, mas bastante subjetiva e permitia uma discussão ampla.
Assim sendo, pode ter dado margem a interpretações diversas e a respostas inexatas, pois demandou uma relação genérica entre os dois conceitos ( gestão ética e transparência ).
O nível de profundidade da questão foi mediano.
O assunto é popular e aparece assiduamente, tanto em discussões acadêmicas, em debates profissionais como nas mídias populares, portanto, bastante atual e não demanda um conhecimento profundo dos conceitos envolvidos.
Contudo, a atualidade do tema é fator que pode levar à banalização da questão, como se percebeu em várias respostas.
Muitos alunos não se preocuparam em elaborar e incorporar conceitos e definições, nem mesmo os elementos apresentados no corpo do enunciado da questão, tais como as características de normas de conduta e valores éticos.
A resposta padrão esperada apresentou uma abordagem lógica : o respondente deveria correlacionar a demanda por uma gestão ética e transparente com a prática da democracia.
Os princípios elencados na questão deveriam servir de liame para esta lógica.
O respondente deveria, portanto, valorizar a ideia de ética e de transparência, particularmente a segunda, enquanto forma de viabilizar a primeira, e, por fim, para completar a resposta, relacionar transparência e ética para a consecução de uma administração pública construída sobre princípios republicanos e democráticos.
Uma elaboração mais preciosa e completa foi rara.
Ademais, o critério de distribuição de pontos não permitiu uma variação maior na nota, pois caso a resposta não apreciasse os primeiros elementos elencados no padrão de resposta, não era possível atribuir qualquer pontuação para os elementos supervenientes e complementares.
Cabe comentar que conceitos e princípios foram fornecidos no enunciado da questão e se os alunos se referissem a eles deveriam ter elementos suficientes para uma resposta satisfatória.
Contudo, muitos estudantes não foram capazes de responder de forma completa e abordar a relação entre gestão ética e transparência, mas, somente, de um ou outro aspecto, isoladamente.
O tempo e o espaço fornecidos para responder à questão foram suficientes, uma vez verificado o padrão de resposta.
Com respeito à diversidade das respostas encontradas, não se identificou dificuldade de compreensão do enunciado.
Não houve muito diversidade nas respostas encontradas.
Contudo, a maioria dos estudantes demonstrou que pode discutir o assunto de forma simplista, mas não consegue elaborar um raciocínio mais aprofundado e fazer ilações e relações de dois conceitos centrais para o debate contemporâneo sobre Gestão Pública.
A linguagem adotada pelos alunos foi adequada do ponto de vista dos termos específicos da Área.
No entanto, muitas respostas parafrasearam o que constava do enunciado, particularmente os princípios e características ali elencados, contudo, não explicaram ou elaboraram respostas completas que contemplassem o significado de cada elemento aportado na parte inicial de forma correlacionada.
As notas dos estudantes foram distribuídas em três níveis : fracas ; medianas e boas.
A abordagem, para responder à questão, mais frequentemente encontrada, foi a repetição de parte do enunciado, ou inclusão de algum elemento mencionado na apresentação do problema.
O principal erro foi não elaborar sobre a relação entre gestão pública e transparência e como uma gestão ética usa instrumentos para incentivar a transparência.
E, reversamente, como a transparência conduz a mais ética na gestão pública, dada a maior visibilidade dos atos do poder público.
Houve muita diversidade na apresentação das soluções, na medida em que muitos respondentes discorreram sobre exemplos, casos de falta de ética ou transparência e deram exemplos diferentes.
Muitos mencionaram o problema da corrupção e deram como exemplo o caso da Petrobrás.
Também houve diversidade quando endereçavam a gestão ética ou a transparência e, mais raramente, na forma como os dois conceitos foram conectados.
Nas respostas consideradas fracas, alguns não mencionaram a questão da ética ou da transparência de forma coerente ou alusiva à questão, mas sim a trataram de forma irreverente e depreciativa.
Outros apenas mencionaram a questão ética, ou a transparência ou os princípios éticos.
Nas respostas consideradas medianas ocorreram alguns cenários.
Alguns trataram da gestão ética e da administração, mas não abordaram a transparência.
Alguns discorrem sobre da gestão ética e da administração, mas não correlacionam os conceitos.
Outros trataram da gestão da transparência e da administração, mas não abordaram a ética.
E, alguns discorrem sobre a gestão ética, mencionam a transparência, mas não correlacionam os conceitos.
As respostas consideradas boas foram apresentadas de maneira completa, pois, nessas, os concluintes discorreram sobre a ética e a transparência, bem como sobre a relação entre ambos os conceitos.
A deficiência mais relevante na formação dos alunos em relação ao conteúdo exigido e à capacidade de expressão desses conhecimentos referiu -se à capacidade de expressão clara e inequívoca e à elaboração de correlação entre dois conceitos centrais ao estudo de Administração ou Gestão Pública, baseadas nos princípios de Administração Pública.
Os alunos também demonstraram ser incapazes de discorrer, conceituar ou valorizar de forma adequada, quer seja o conceito de gestão ética ou o de transparência, quer seja sua importância e centralidade para promover"o ethos republicano e democrático como norteador de uma formação que \textbf{ultrapasse} a ética profissional, remetendo -se à responsabilidade pela res pública e à defesa do efetivo caráter público e democrático do Estado"( Diretrizes Curriculares Nacionais do curso de graduação em Administração Pública, de 13/01/2014 ).
Embora este seja um princípio fundamental a ser contemplado em cursos de graduação em Administração Pública, serve igualmente de referência para os cursos de Tecnologia em Gestão Pública.
Quando discorrem e valorizam tais conceitos, como ética e transparência, muitos alunos não são capazes de elaborar nem mesmo sobre elementos trazidos na própria questão, como normas de conduta ou sobre valores éticos, e como eles informam a gestão pública.
Não foi constatada a existência de relação entre esta questão discursiva e alguma questão objetiva da prova.
Em torno de 20 \% dos documentos foram deixados em branco ou foram desconsiderados.
Essa foi a questão com o maior volume de documentos corrigidos. 3.3.2.5 Análise da Questão Discursiva 5 do Componente de Conhecimento Específico A Tabela 3.20 contém as informações relativas à questão 5 do conjunto do Componente de Conhecimento Específico.
A nota Média dos estudantes de todo o Brasil foi 41,0.
A maior Média foi registrada na região Norte ( 51,9 ), enquanto a menor Média foi registrada na região Sul ( 38,2 ).
Quanto à variabilidade das notas, o Desvio padrão dos alunos do Brasil, como um todo, foi 35,9.
Enquanto o maior desvio foi encontrado na região Sudeste ( 36,9 ), o menor foi encontrado na região Norte ( 35,0 ).
As notas máximas ( 100,0 ) e as notas mínimas ( 0,0 ) foram iguais para todas as regiões, sem exceção.
A Mediana para o Brasil com um todo foi 30,0, a mesma obtida nas regiões Sul e Centro-Oeste, enquanto nas demais regiões, a Mediana foi 60,0.
O Gráfico 3.13 representa a distribuição das notas da questão discursiva 5 do Componente de Conhecimento Específico.
Destacam -se o grande número de estudantes que deixou esta questão em branco, mais de 22,5 \% do total, e o intervalo ( 20 ; 30 ], que corresponde à moda da distribuição. 3.3.2.6 Comentários sobre as respostas à Questão Discursiva 5 O comando da questão solicitava um texto dissertativo que versasse sobre a relação da participação social no processo de planejamento governamental e a importância de fatores culturais.
Este conteúdo é assaz pertinente ao estudo de Gestão Pública.
Elaborar sobre o assunto, ser capaz de relacionar a importância da participação social e da inclusão de fatores culturais no programa político e no Plano Plurianual, entender a correlação dos dois conceitos são conteúdos de \textbf{aprendizado} obrigatório no estudo de Gestão Pública.
São elementos basilares no ensino de Gestão Pública contemporânea e parte das competências e habilidades elencadas nas Diretrizes Curriculares da Área de Administração Pública.
O conteúdo deveria ser fundamental e básico nos cursos de Administração e Gestão Pública, particularmente estabelecido no Catálogo Nacional dos Cursos Superiores de Tecnologia para os cursos de Tecnologia em Gestão Pública ( Portaria Inep nº 231, de 10 de junho de 2015, publicada no Diário Oficial de 12 de junho de 2015, Seção 1, pág. 22 ).
O nível de profundidade da questão foi de mediano a difícil.
Apesar de o tema ser recorrentemente associado à cultura e à participação social, em um exame específico da Área é preciso evitar -se a banalização, como ocorreu em várias respostas.
O próprio enunciado permitia uma abordagem genérica do termo cultura, por exemplo.
Nesta questão, os alunos não incorporaram conceitos e definições e, quando usaram exemplos, o fizeram de forma equivocada.
Por exemplo, citaram projetos culturais como se tais fossem a formação e o arcabouço cultural per se.
Os respondentes também mencionaram agricultura familiar de forma errática, sem trazer aspectos culturais ou de participação social para justificar a inclusão no planejamento governamental.
A resposta padrão sugeria que o estudante deveria discorrer sobre a importância da participação social para contemplar no plano de governo as necessidades da sociedade ( no caso, o PPA, processo de planejamento governamental ) e enfatizar a importância do processo participativo como forma de \textbf{aprendizado} para gestores públicos se informarem sobre demandas e necessidades da população.
Deveriam, portanto, discorrer sobre a valorização de valores e aspectos culturais, como fator importante para sinalizar a existência de grupos diversos e demandas específicas e distintas necessidades.
Por sua vez, também se fazia necessário indicar que, prestar atenção às questões culturais pode influenciar o sucesso ou insucesso do planejamento.
O respondente deveria, ainda, destacar a ideia de participação social como processo de aprendizagem e informação para os agentes públicos e políticos e, por fim, para completar a resposta, elaborar sobre cultura como fator social diferenciador entre grupos sociais e determinantes do sucesso de uma política.
O critério de distribuição da pontuação forneceu uma variação gradual e incremental na nota, que permitiu avaliar cada item da questão separadamente.
Como foram avaliações paralelas, as respostas de cada item foram analisadas de forma independente, com pontuação atribuída em separado, beneficiando o respondente de forma mais justa, apesar de avaliar o conhecimento de forma isolada ou estanque, e não interativa e dinâmica, em uma relação de \textbf{causalidade} com o processo democrático e republicano.
Cabe ressaltar que os exemplos foram fornecidos no enunciado da questão e teria sido suficiente relacioná -los de forma correta para responder de forma satisfatória.
Contudo, muitos não foram capazes de responder de forma relacional e mencionar a conexão entre participação social, sua importância para definir projetos e informar gestores públicos, e, ainda, acrescentar a importância dos fatores culturais.
A maioria simplesmente abordou um ou outro aspecto, isoladamente.
O tempo e o espaço fornecidos para responder à questão foram suficientes e de acordo com o padrão de resposta.
De acordo com o padrão de resposta, os critérios utilizados para a correção deveriam consistir em observar a presença dos itens a serem indicados para pontuação na resposta padrão.
No Item ‘ a ’, as respostas deveriam discorrer, mesmo que sucintamente, sobre aspectos, tais como : a ) a participação social como elemento essencial ao aprimoramento das políticas públicas ; b ) a participação social conduz ao exercício da cidadania e ao fortalecimento da democracia na gestão pública, e c ) a participação social como instrumento positivo para a definição de prioridades das ações governamentais.
No item ‘ b ’, o respondente deveria discorrer sobre os fatores culturais como elementos reveladores de características específicas e particulares de grupos sociais a serem incluídos na elaboração de políticas públicas, como informadores de demandas e necessidades específicas de um grupo social, local ou regional.
As respostas deveriam considerar que as diferenças culturais devem ser levadas em consideração para o sucesso das políticas públicas.
Ademais, os fatores culturais permitiriam não somente evidenciar o modo de vida, as necessidades e peculiaridades locais, mas também revelar oportunidades que podem ser incluídas no planejamento governamental.
Apesar de ter sido estabelecido um critério de distribuição de pontos mais flexível, ainda assim houve dificuldade em pontuar as questões em uma escala com tão poucos “ degraus ”.
A orientação permitiu considerar uma resposta que mencionasse, de forma simplista, os valores culturais com a nota máxima.
Isto diminuiu a chance de graduar nuances e contribuições significativas, mas consideradas de menor valor.
A linguagem adotada pelos alunos foi adequada, do ponto de vista dos termos específicos da Área.
As notas dos estudantes foram distribuídas em três níveis : fracas ; medianas e boas.
Considerando -se que quase metade dos respondentes teve um desempenho fraco nesta questão, deve -se avaliar cuidadosamente o enunciado, o padrão de resposta e os critérios de atribuição de pontos estabelecidos, já que possivelmente contribuíram para tal resultado.
Evidentemente, deve -se considerar o comentário acima em conjunção com o despreparo dos alunos e a incapacidade de articular conceitos e relações de causa de efeito.
A abordagem mais frequente foi enfatizar a importância da participação social.
O principal erro encontrado foi confundir eventos culturais com cultura e valores culturais per se.
Muitos mencionaram turistas e agricultura familiar como projetos culturais.
Outros mencionaram festejos, eventos musicais e danças como cultura, e não como elementos de cultura específica de uma região.
Além disso, não houve muita diversidade nas soluções apresentadas.
Dentre as respostas consideradas fracas, as que receberam nota zero os respondentes não mencionaram participação social ou cultura.
Alguns respondentes mencionaram aspectos relevantes sobre participação social.
Nas respostas consideradas medianas, os respondentes elaboraram resposta somente sobre participação social e sua importância para influenciar gestores e políticos na elaboração do planejamento governamental.
Não fizeram referência aos valores culturais Dentre as respostas consideradas boas, alguns respondentes abordaram a participação social sem mencionar o processo de influenciar os dirigentes políticos e mencionaram a importância dos valores culturais.
E as respostas completas, que receberam o total de pontos, incluíam elaboração sobre participação social e sobre aspectos culturais influenciando o planejamento governamental.
As respostas revelam que os alunos são incapazes de articular vários aspectos como vetores do sucesso ou insucesso de uma política pública.
A maioria consegue elaborar uma resposta sobre somente um aspecto.
Também quanto à capacidade de expressão desses conhecimentos, os alunos mostraram -se pouco preparados.
Eles não conseguem elaborar um texto mais complexo e sofisticado.
A questão era interessante e pertinente, mas foi mal formulada e pode ter levado o aluno a concluir que eventos culturais e cultura são sinônimos.
Fosse a questão melhor estruturada, talvez os alunos tivessem se apropriado do enunciado de forma mais correta. 3.3.3 Considerações Finais A Portaria Inep nº 231, de 10 de junho de 2015, publicada no Diário Oficial de 12 de junho de 2015, Seção 1, pág. 22, destaca em seu artigo 4º que : “ A prova do Enade/2015, no Componente Específico da Área de Tecnologia em Gestão Pública, terá como subsídio a descrição do perfil do tecnólogo oferecida pelo Catálogo Nacional dos Cursos Superiores de Tecnologia para os cursos de Tecnologia em Gestão Pública, bem como normativas associadas e a legislação profissional ”.
A seguir, nos artigos 5º a 7º, a Portaria enuncia os componentes específicos da área ( artigo 5º ), as competências e habilidades ( artigo 6º ) e os componentes específicos e os conteúdos curriculares ( artigo 7º ).
A seguir, faz -se algumas considerações sobre as questões propostas no ENADE/2015, para a área de Tecnologia em Gestão Pública, vis-à-vis os preceitos normativos estabelecidos pelo Instituto Nacional de Estudos e Pesquisas Educacionais Anísio Teixeira ( Inep ).
Quanto aos componentes específicos da Área, as questões elaboradas, incorporaram aspectos elencados no item I, como a análise do cenário político ( Questões 4 e 5 ) e o econômico ( Questão 3 ) e a compreensão sistêmica ( item II, do artigo ), aplicável em todas as questões.
O resultado final demonstrou que os programas de Tecnologia em Gestão Pública devem rever e enfatizar tais aspectos em seus currículos, de forma a levar os profissionais a"tomarem decisões em um contexto diversificado"e a promoverem o"entendimento da complexa dimensão de suas atribuições perante o Estado e a sociedade, tornando -se eficientes e éticos no exercício de suas funções, no gerenciamento das relações sociais e dos recursos humanos, administrativos e de serviços nas organizações públicas".
No quesito sobre competências e habilidades, listado no artigo 6º, itens I a X, da citada Portaria, faz -se evidente a necessidade de promover maior efetividade na comunicação e expressão escrita ( item I ) ; incentivar o uso da lógica e coerência ao citar os conceitos referentes aos fundamentos teóricos da Administração ( item II ) ; relacionar melhor as respostas com o caráter social das questões político-econômicas ( item V ) ; valorizar a elaboração de políticas públicas por meio de planos, programas e projetos ( VI ), privilegiando as demandas da sociedade ( item VII ) ; e relacionar a consciência quanto às implicações éticas do exercício profissional, em especial a compreensão do ethos republicano e democrático, indispensável à sua atuação ( item VIII ).
Estes elementos devem estar presentes em todas os programas e os alunos deveriam ser capazes de expressá -los"de forma assertiva, efetiva e legítima"( item X ).
De uma forma geral, poucos foram capazes de demonstrar que conseguiram apreender, atingir e desenvolver as competências e habilidades almejadas de forma completa para um concluinte de um programa de Tecnologia em Gestão Pública.
A compreensão sistêmica e integrada dos conceitos, assim como a capacidade de articular de forma efetiva e assertiva os conceitos e as correlações no campo de Gestão Pública não são competências e habilidades atingidas pela grande maioria dos alunos.
Finalmente, no que diz respeito ao conteúdo curricular, só foi possível analisar os quesitos sobre comunicação e linguagem ; governo e administração pública, políticas públicas, economia do setor público ; e ética ( incisos I, VII, XIII, XVI, XVII, respectivamente ).
A análise ficou restrita às questões elaboradas e, portanto, abordou somente alguns elementos dos estudos na área, sendo muitas vezes tangencial.
Conceitos e teorias foram pouco explorados nas respostas.
Contudo, o fato de o conjunto das questões contemplar uma questão de economia do setor público, uma sobre ética e planejamento e uma terceira sobre participação social e cultura, mostra a amplitude dos conhecimentos exigidos na Área e cobriu conteúdos curriculares apropriados.
Assim, pode -se dizer que quanto a este quesito, a prova examinou efetivamente o conhecimento específico da Área.
Contudo, quanto à mensuração deste conhecimento, a comunicação e linguagem utilizadas pelos alunos, pode -se dizer que os resultados indicaram um desempenho pior e menos efetivo do que deveria ter sido alcançado.
Acredita -se que se os programas e cursos de Tecnologia em Gestão Pública incluíssem a divulgação dos propósitos, objetivos de aprendizagem e as competências estabelecidas na Portaria aqui mencionada, de forma enfática e recorrente, os alunos teriam sido capazes de estabelecer um raciocínio lógico mais preciso.
Também, ênfase deve ser colocada no exercício contínuo de uma comunicação efetiva, assertiva e competente, uma vez que emerge das respostas uma incapacidade genérica de estabelecer relações complexas e comunicá -las de forma compreensível e sistemicamente organizada.

% matched lemmas: aprendizado, causalidade, desordem, exacerbar, fraternidade, intérprete, liberal, morfema, nocivo, novembro, peso, ultrapassar
\end{textsample}
